% !TEX TS-program = xelatex
\documentclass[10pt, letterpaper]{article}

% -------------------------
% Page / Layout
% -------------------------
\usepackage[
ignoreheadfoot,
top=1.5cm,
bottom=1.7cm,
left=1.7cm,
right=1.7cm,
footskip=1.0cm
]{geometry}

\usepackage{tabularx}
\usepackage{array}
\usepackage{titlesec}
\usepackage{enumitem}
\usepackage{graphicx}
\usepackage{calc}
\usepackage[dvipsnames]{xcolor}
\definecolor{primaryColor}{RGB}{0,0,0}
\usepackage[colorlinks=true, urlcolor=primaryColor]{hyperref}

% -------------------------
% Korean (IMPORTANT)
% -------------------------
% Use XeLaTeX/LuaLaTeX + kotex for robust Korean typesetting
\usepackage{kotex} % loads xetexko/luatexko as needed
% (Optionally set fonts; comment out if you want system defaults)
% \setmainhangulfont{NanumMyeongjo}
% \setsanshangulfont{NanumGothic}
\usepackage{changepage} % provides adjustwidth
\usepackage{tikz}

\newenvironment{onecolentry}{
	\begin{adjustwidth}{0cm}{0cm}
	}{
	\end{adjustwidth}
}


% -------------------------
% Style
% -------------------------
\pagestyle{empty}
\setcounter{secnumdepth}{0}
\setlength{\parindent}{0pt}
\setlength{\columnsep}{0.15cm}

\titleformat{\section}{\bfseries\large}{}{0pt}{}[\vspace{2pt}\titlerule]
\titlespacing{\section}{0pt}{0.35cm}{0.15cm}

\renewcommand\labelitemi{$\vcenter{\hbox{\small$\bullet$}}$}
\newenvironment{highlights}{
	\begin{itemize}[
		topsep=0.10cm,
		parsep=0.05cm,
		partopsep=0pt,
		itemsep=0pt,
		leftmargin=14pt
		]
	}{
	\end{itemize}
}

% A compact row style for Korean resume tables
\renewcommand{\arraystretch}{1.25}
\newcolumntype{L}[1]{>{\raggedright\arraybackslash}p{#1}}
\newcolumntype{Y}{>{\raggedright\arraybackslash}X}

% -------------------------
% Document
% -------------------------
\begin{document}
	
% =========================
% TOP HEADER (CV style): Name + tagline + inline contacts + PHOTO
% =========================
\noindent
\begin{minipage}[t]{\dimexpr\textwidth-3.8cm\relax}
	\vspace{0pt}
	
	{\fontsize{18pt}{18pt}\selectfont\textbf{Swain Bishal Ranjan}\par}
	\vspace{2pt}
	{\small Kumoh National Institute of Technology (KIT) · 박사\par}
	\vspace{6pt}
	
	{\small
		\href{mailto:bishalswain@kumoh.ac.kr}{bishalswain@kumoh.ac.kr}\quad|\quad
		010-4683-3356\quad|\quad
		경북 구미시 옥계동 801-7\par
		\href{https://www.bishalswa.in/}{www.bishalswa.in}\quad|\quad
		\href{https://github.com/bluesaiyancodes}{github.com/bluesaiyancodes}\quad|\quad
		\href{https://www.linkedin.com/in/bishal-swain/}{linkedin.com/in/bishal-swain}\par
		\textbf{국적:} 인도\par
	}
\end{minipage}\hfill
\begin{minipage}[t]{4.0cm}
	\vspace{0pt}
	\centering
	\begin{tikzpicture}
		% circle radius
		\def\r{1.55cm}
		
		% clip and shift image slightly upward for better face framing
		\clip (0,0) circle (\r);
		\node at (0,0.10) {\includegraphics[width=3.3cm]{photo.jpg}}; % adjust 0.10 up/down
		
		% thin border (draw after clip)
		%\draw[line width=0.4pt] (0,0) circle (\r);
	\end{tikzpicture}
\end{minipage}


\vspace{-0.5cm}

% =========================
% SUMMARY (Korean recruiter-friendly)
% =========================

\section{요약}
\begin{highlights}
	\item 컴퓨터비전 박사과정 연구자: 분할(Segmentation) 및 구조적 예측을 중심으로, 도메인 변화(OOD)와 지속적 학습(Continual Adaptation)에 강건한 모델 설계 연구
	\item 연속시간 동역학 기반 상태 진화와 연관 메모리 메커니즘을 결합한 아키텍처 탐구 — 입력 의존적 시간적 적응과 메모리 피드백을 통한 일반화 및 안정성 향상
	\item Foundation Model·Vision-Language Model 기반 분할 및 멀티스케일 표현 학습, 이론적 안정성 분석과 벤치마크 기반 실험 검증 병행
	\item MICCAI Student Board Social Vice-President 활동 %; 정주초청외국인 장학생 (GKS)
\end{highlights}

	
	% =========================
	% EDUCATION (학력)
	% =========================
	\section{학력}
	\begin{tabularx}{\textwidth}{@{}L{3.2cm}Y L{4.5cm}@{}}
		\textbf{기간} & \textbf{학교/전공} & \textbf{비고} \\
		2022.09 -- 2026.08 & 국립금오공과대학교 / 컴퓨터$\cdot$AI융합공학과 & 지도교수: 고재필 \\
		2021.09 -- 2022.08 & 전남대학교 / 언어교육원 & 한국어(4 단계) \\
		2018.07 -- 2020.05 & 인도국립 폰디체리대학교 / 컴퓨터과학 석사(M.Sc.) & GPA 9.36/10.0 (우등 졸업) \\
	\end{tabularx}
	
	% =========================
	% EXPERIENCE (경력/연구경력)
	% =========================
	\section{경력/연구경력}
	\begin{tabularx}{\textwidth}{@{}L{3.2cm}Y L{4.5cm}@{}}
		\textbf{기간} & \textbf{기관/직무} & \textbf{주요내용} \\
		2020.06 -- 2021.07 & 인도과학대학교 (IISc), CDS / Research Associate & ML 연구 및 엔지니어링 협업 \\
		2019.11 -- 2020.05 & 인도우주연구기구(ISRO) 산하 국립대기연구소(NARL) / Research Intern & 대기·풍장 데이터 분석 및 시각화 연구 \\
	\end{tabularx}
	

	
	% =========================
	% PUBLICATIONS (논문/발표)
	% =========================
	\section{최근 논문}
	\begin{highlights}
		\item Memory-Induced Attractor Dynamics in Liquid Time-Constant Networks, \textit{ICANN 2026}. (학회)
		\item Coupling Liquid Time-Constant Encoders with Modern Hopfield Memory, \textit{\textbf{CVPR 2026}}. (학회)
		\item Guided Texture Segmentation via Coordinate-Aware Class-Ratio Mapping, \textit{\textbf{WACV 2026}}. (학회)
		\item Out-of-Distribution Nuclei Segmentation in Histology Imaging via Liquid Neural Networks with Modern Hopfield Layer, \textit{\textbf{MICCAI 2025}}. (학회)
		\item SAM Guided Task-Specific Enhanced Nuclei Segmentation in Digital Pathology, \textit{\textbf{MICCAI 2024}}. (학회)
		\item Nuclei Segmentation in Histopathological Images with Enhanced U-Net3+, \textit{\textbf{MIDL 2024}}. (학회)
		%\item Phase Fraction Calculation in Multiphase Steel Microstructures Through Semantic Segmentation, \textit{ICASI 2024}.
		\item Complex-Phase Steel Microstructure Segmentation Using UNet: Analysis across Different Magnifications and Steel Types, \textit{MDPI Materials 2023}. (SCIE급 논문)
		%\item Rise of Fluid Computing: A Collective Effort Of Mist, Fog and Cloud, \textit{IJCSE 2019}.
	\end{highlights}
	
	% =========================
	% PROJECTS (프로젝트)
	% =========================
	%\section{프로젝트}
	%\begin{highlights}
	%	\item \textbf{Alloy Steel Microstructure Segmentation (PL)}: CCRM + Adaptive Gate Fusion, 평균 Dice +6.1\%, 파라미터 증가 <2\%
	%	\item \textbf{Pipe Defect Detection (PL)}: SAM + LoRA + box prompt 기반 내부 결함 분할, 이미지 단위 Precision 1.00 / Recall 0.94
	%	\item \textbf{Battery Particle Segmentation (PL)}: SAM 반자동 라벨 + LoRA 적응, 평균 test Dice 96.14
	%	\item \textbf{Temperature Super-Resolution (PL)}: Residual U-Net + PixelShuffle 4$\times$ 업스케일링, test L2 0.0139
	%	\item \textbf{ClimaX via Radar}: RegionalClimaX 기반 레이더 강수 예측(최대 72h), 테스트 정확도 62.1\% (multi-var)
	%\end{highlights}
	
	% =========================
	% CERTS / LANGUAGE (자격/어학)
	% =========================
	\section{수상/자격}
	\begin{highlights}
		\item 2025학년도 후기 대학원 졸업 포상 - 우수논문상 (2026)
		\item Global Korea Scholarship (GKS) - 정주초청외국인 장학생 (2021-2026)
		\item TOPIK-5 (2025)
		\item KIIP 사회통합평가 - 4급 (2025)
		\item Docker \& Kubernetes (Udemy, 2020), DeepLearning.AI Specializations (2019--2020)
		\item RHCSA (2016), EC-Council Certifications (2017--2018)
	\end{highlights}
	
	% =========================
	% SKILLS (보유기술)
	% =========================
	\section{보유기술}
	\begin{onecolentry}
		\textbf{언어/프레임워크:} Python, PyTorch, TensorFlow \\
		\textbf{CV/ML}: OOD Segmentation, Vision-Language Models, Continuous-time Neural Networks, Continual Learning \\
		%\textbf{개발환경:} Linux, Git, CUDA, Docker, W\&B, 실험관리 및 재현성 중심 연구 파이프라인 구축
	\end{onecolentry}
	
	% =========================
	% Voluntary / Others
	% =========================
	%\section{기타}
	
	
\end{document}
